\documentclass[11pt]{article}
\usepackage[a4paper,margin=24mm]{geometry}
\usepackage{amsmath,amssymb,amsthm,hyperref}
\newtheorem{theorem}{Theorem}
\newtheorem{lemma}{Lemma}
\newcommand{\PP}{\mathcal P}
\title{Synchronized local completion with a fixed number of controlled moments}
\author{Complete proof; three independent mathematical reviews completed}
\date{30 September 2026}
\begin{document}\maketitle
Fix $q\ge1$ and positive constants $a,b$. Polynomial bounds below are
in one parameter $P$, with exponents depending only on $q,a,b$ and the
fixed input polynomial bounds. A single common denominator, rather
than separate bounds on individual denominators, is part of the conclusion.

\begin{theorem}\label{main}
Consider polynomially many rational polynomial probability densities
$w_\alpha$ on $[0,1]$, satisfying $a\le w_\alpha\le b$, with polynomial
degrees, and let $m\ge\max(2,q)$ be polynomially bounded.
Put $R=m+q+2+\max_\alpha\deg w_\alpha$.
For each instance prescribe an integer $K_\alpha$ with
\[
 CR^4\le K_\alpha\le\operatorname{poly}(P).
\]
Suppose the numbers
\begin{equation}\label{inputarithmetic}
 K_\alpha\int_0^1t^j w_\alpha(t)dt,
       \qquad 0\le j\le2q+1,
\end{equation}
have one polynomial common denominator. There exist one polynomial
integer $L$ and, for every instance, rational pure densities
$f_{A,\alpha},f_{B,\alpha}$ such that:
\begin{enumerate}
\item Their ordered polynomial moments through degree $m$ equal those
of independent uniform and $w_\alpha$ variables, respectively.
\item $B$ occupies $K_\alpha$ disjoint boxes, each of mass $1/K_\alpha$;
$A$ occupies the complementary gaps. The endpoints of \emph{all}
boxes and gaps across \emph{all} instances lie on one polynomial grid.
\item On every $A$ gap, all unnormalized moments through $q$ are
integers divided by $L$. On every $B$ box, all conditional moments
through $q$ have one polynomial common denominator, shared across all
boxes and all instances.
\item The numbers of cells, reciprocal lengths, reciprocal positive
masses, and polynomial degrees are polynomially bounded. Conditional
densities, rescaled to their cells, have fixed positive upper and lower
bounds depending only on $a,b,q$.
\end{enumerate}
No small-denominator claim is made for higher moments or coefficients.
\end{theorem}

\begin{lemma}[A fixed finite set of forced quadrature nodes]\label{forced}
For each fixed $h$, one may strengthen the separated weighted quadrature
lemma in \path{generalized_local_two_variable_completion.tex} by
requiring nodes at the prescribed positions
$\tau_i=i/(h+1)$, $1\le i\le h$.
Every sufficiently large $K\ge C_h(D+d+1)^4$ is allowed for a degree-$D$
rule of a degree-$d$ density with bounds $a,b$.
All endpoint and consecutive gaps remain at least $c_h/K$.
\end{lemma}
\begin{proof}
Begin with the midpoint quantile grid. For each $\tau_i$, replace a
nearest grid point by $\tau_i$; these chosen points are distinct once
$K$ is large. A nearest-point choice leaves distance at least half the
minimum original grid spacing to every unchosen neighbor. Each move
is $O_{a,b}(1/K)$. Its contribution to the normalized moment residual
is $O(\operatorname{poly}(D+d)/K^2)$, and there are only $h$ moves.
Run the same fixed-right-inverse Newton argument on the \emph{other}
nodes, holding these $h$ coordinates fixed. Deleting $h$ columns from
the ordinary Legendre derivative Gram changes it by at most $C_hD^2/K$;
the forced moves produce another polynomially small error. The lower
Gram bound survives. The correction radius is still
$C_h(D+d+1)^4/K^2$, smaller than a prescribed small multiple of $1/K$.
Thus the free nodes remain separated, and the fixed nodes are retained.
\end{proof}

\begin{proof}[Proof of Theorem~\ref{main}]
\emph{One universal feature family.}
Use the cumulative-moment conditioning theorem in
\path{../rational_designs/asymptotic/multiple_prefix_moment_conditioning.tex}.
Choose the same $q$ separated Legendre bands for every instance, with
the first degree larger than $\max\deg w_\alpha+J+q+1$, where
$J=m+q+1$. All band degrees are $O_q(R)$.
With one common sufficiently small inverse-polynomial $\eta$, put
\[
 f_*(t)=1-\eta\sum_{i=1}^qP_{d_i}(2t-1),\qquad
 M_s(t)=\int_0^t u^s f_*(u)du\quad(0\le s\le q).
\]
These features are within $C_q\eta$ of $t^{s+1}/(s+1)$ and have
polynomially bounded derivatives. Their polynomial-multiplier family
$1,M_0,\ldots,M_q$, with multiplier degrees at most $J-1$, has an
inverse-polynomial Gram lower bound. This holds uniformly with every
weight $w_\alpha$ by the assumed density bounds.
Legendre orthogonality gives
\begin{equation}\label{compatbaseline}
 \int t^j M_s(t)w_\alpha(t)dt
  =\frac1{s+1}\int t^{j+s+1}w_\alpha(t)dt
       \qquad(0\le j\le J,\ 0\le s\le q).
\end{equation}

Set $h=q+2$. Use Lemma~\ref{forced} to choose, for each prescribed
$K_\alpha$, a separated real rule $b^0_{\alpha u}$ of sufficiently high
degree $O_q(R)$. It is exact for every moment in
\eqref{compatbaseline} and for squares of every feature-multiplier
polynomial just described. It contains all $h$ universal control nodes
$\tau_i=i/(h+1)$.

\emph{Shared rational control features and exact rank sums.}
Let $D_0$ clear the common denominators in \eqref{inputarithmetic}.
Choose one large polynomial integer $L$ divisible by the PRODUCT
$D_0(q+1)!$. At the $h$
control indices, round the numbers $M_s(\tau_i)$ DOWN once to multiples
of $1/L$, using identical values in every instance. These feature
values, together with the constant feature $1$, form a square
$h\times h$ matrix $\Psi_*$. It is a small perturbation of the fixed
Vandermonde-type matrix with rows
$1,\tau,\tau^2/2,\ldots,\tau^{q+1}/(q+1)$.
Hence its inverse norm is bounded in terms of $q$ alone. Its rational
inverse has one polynomial denominator, because $h$ is fixed and all
entries have denominator dividing $L$.

Round the other feature values $M_s(b^0_{\alpha u})$ to the same grid,
and correct their sums to satisfy
\[
 \sum_u r_{s,\alpha u}
 =\frac{K_\alpha}{s+1}\int t^{s+1}w_\alpha(t)dt.
\]
The desired sums multiplied by $L$ are integers by
\eqref{inputarithmetic}. The control values remain fixed. After
rounding down, at most $K_\alpha$ integer units are needed, so adding
at most two units to each of the $K_\alpha-h$ free entries suffices
for large enough $K_\alpha$. Thus every feature changes by at most
$2/L$. The same assertion holds for each of the fixed number of
feature rows independently.

Hold all these rational features fixed and correct the nodes to
restore every nonconstant ordinary and feature-weighted moment through
$J$. In the antiderivative Legendre basis, the derivative evaluates
$1,M_0,\ldots,M_q$ times the Legendre polynomials of degrees below $J$.
Its baseline empirical Gram is exactly the weighted continuous Gram.
The conditioning theorem, polynomial derivative bounds, and the
$O(\operatorname{poly}(P)/L)$ rounding residual give the same
fixed-right-inverse contraction as in the reviewed local two-variable
proof. With a sufficiently large exponent for $L$, the node displacement
is smaller than any prescribed inverse polynomial. All separations
survive. The control \emph{nodes} may move at this step; their rational
\emph{features} remain the identical columns of $\Psi_*$, which is
what the denominator argument below requires.
We now have exact ordinary moments and
\begin{equation}\label{rankexact}
 K_\alpha^{-1}\sum_u r_{s,\alpha u}b_{\alpha u}^j
 =\frac1{s+1}\int t^{j+s+1}w_\alpha(t)dt,
       \qquad 0\le j\le J.
\end{equation}
All feature arrays remain within $C_q\eta$ of the uniform prefix
moments $b^{s+1}/(s+1)$, after choosing the rounding precision finer.

\emph{Smoothing with one common low-moment denominator.}
Apply the fixed-$q$ box-smoothing argument in
\path{../rational_designs/asymptotic/polynomial_mean_box_smoothing.tex}
using the feature rows $(1,r_0,\ldots,r_q)$.
Choose one common inverse-polynomial halfwidth $\epsilon$ and one
common polynomial center-rounding grid, sufficiently fine for all
instances. Uniform bounds on all preceding condition numbers allow
these choices. First solve the real uniform-box moment problem;
then round centers and low conditional box moments through $q$ on
common grids, and correct the low global moment totals using exactly
the control minor $\Psi_*$ in each instance.

Here is the arithmetic point in full. The control equations have form
\[
 \Psi_*\mu_{\rm control,j}
 =K_\alpha T_{\alpha,j}
       -\sum_{u\notin\mathrm{control}}\psi_{\alpha u}\mu_{\alpha u,j}.
\]
Every entry on the right has one common polynomial denominator across
all instances. For the sum this follows from the common feature and
moment grids. For the targets it follows from
\eqref{inputarithmetic}, since ordinary targets are
$K_\alpha\int t^j w_\alpha$ and feature targets are
$K_\alpha\int t^{j+s+1}w_\alpha/(s+1)$, with $j\le q$ and
$j+s+1\le2q+1$.
The inverse of $\Psi_*$ is the SAME rational matrix in every instance.
Therefore all corrected conditional low box moments share one
polynomial denominator. No least common multiple of unrelated minors
is used. The remaining high moments are corrected by the degree-$(q+1)$
Legendre mode on each box; its rational coefficients may have unrestricted
denominators. The reviewed argument gives positive conditional densities
with fixed bounds, exact ordinary moments and exact \eqref{rankexact}.
All box endpoints use the common center grid and halfwidth.

\emph{Positive completion on the gaps.}
Prescribe on the $A$ gaps cumulative moments $r_{s,\alpha u}$,
$0\le s\le q$, before each $B$ box, with the endpoint values
$0$ and $1/(s+1)$. All lie on the $1/L$ grid. On each gap realize
the required moments by a polynomial of degree $q$, obtaining $g$.
The elementary rescaled moment system gives
\[
 \|g-1\|_{L^\infty(\mathrm{gaps})}
 \le \operatorname{poly}_q(P)(\eta+\epsilon+	ext{node/center errors}),
\]
while extending $g$ by zero on the boxes adds at most
$\sqrt{2K_\alpha\epsilon}$ to $\|1-g\|_2$.
The exact equations \eqref{rankexact}, and the ordinary box moments,
are precisely the compatibility conditions on
\[
 \mathcal S_q=\PP_q+\operatorname{span}\{a^sS_t(a):s\le q,\ s+t\le m\}.
\]
Use the fixed-$q$ quantitative smoothed quotient and positive-completion
extension in
\path{../independent_directions/quantitative_smoothed_quotient.tex}.
Its losses are polynomial for fixed $q$ and inverse-polynomial gap
lengths. Choose first $\eta$ sufficiently small, then $\epsilon$,
then all rounding precisions sufficiently fine. The resulting rational
Riesz correction has zero moments through $q$ on every gap and produces
$f_A$ between $1/4$ and $7/4$ there. Thus each gap moment through $q$
is exactly a difference of entries of the $1/L$ prefix grids.
The density is polynomial of degree at most $m$ on gaps; $B$ pieces
have degree at most $q+1$. Lengths and masses are inverse-polynomial,
and conditional density bounds follow by normalizing each cell.

Finally, exact integration of
$\PP_m+\operatorname{span}\{a^rS_s(a):r+s\le m\}$ gives
\[
 \mathbb E[A^rB^s\mathbf1_{A<B}]
 =\frac1{r+1}\int b^{r+s+1}w_\alpha(b)db
       \qquad(r+s\le m),
\]
and the ordinary and opposite-chamber moments follow as before.
This proves the ordered target assertion and all denominator claims.
\end{proof}
\end{document}
